\documentclass[english,twoside, openright, hidelinks, 11 pt, class=report,crop=false]{standalone}
\input{../../preamb}
\input{../bm}

\begin{document}
\opr{opgbakke} \\
\textbf{Alternativ 1} \os
\abc{
\item $\frac{7}{24}\approx\frac{8}{24}=\frac{1}{3}\approx0,33$. Stigningsgraden er ca. $33\%$. \os

\textbf{Alternativ 2} \os
$\frac{7}{24}\approx\frac{7}{21}=\frac{1}{3}=0,33$. Stigningsgraden er ca. $33\%$.\os

\textbf{Alternativ 3}\os
$\frac{7}{24}\approx\frac{6}{24}=\frac{1}{4}=0,25$. Stigningsgraden er ca. $25\%$. \os

\textbf{Alternativ 4}\os
$\frac{7}{24}\approx\frac{6}{24}=\frac{1}{4}=0,25$ og $\frac{7}{24}\approx\frac{8}{24}=\frac{1}{3}\approx0,33$. Stigningsgraden er ca. midtpunktet mellom $ 25\%$ og $33\%$, som er $29\%$.
\item Veien er hypotenusen i figuren. Av Pytagoras' setning er
\[\text{tallverdi til veilengde}= \sqrt{24^2+7^2}=\sqrt{476+49}=\sqrt{625}=25 \]
Altså er veien 25\enh{km}. \os
\textbf{Alternativ 1}\os
Hvis man kjører i 60\enh{km/h}, bruker man 1\enh{min} per kilometer. Dette betyr 25\enh{min} på 25\enh{km}. Siden bilen brukte 24\enh{min}, må den ha kjørt raskere enn 60\enh{km/h}, og dermed brutt fartsgrensen. \os

\textbf{Alternativ 2} \os
Vi har at $24\enh{min}=\frac{24}{60}\enh{h}=\frac{2}{5}\enh{h}$. Videre er
\[ \text{fart}=\frac{\text{lengde}}{\text{tid}}=\frac{25}{\frac{2}{5}}=25\cdot\frac{5}{2}=62,5 \]
}
Farten var altså $62,5\enh{h}$, og dermed brøt bilen fartsgrensen.\os

\textbf{Alternativ 3}\os
Vi har at $24\enh{min}=\frac{24}{60}\enh{h}=\frac{2}{5}\enh{h}$.
Hvis bilen holdt fartsgrensen ville lengden være gitt som
\[ \text{lengde}=\text{fart}\cdot\text{tid}=50\cdot\frac{2}{5}=20 \]
Ved å holde fartsgrensen vill altså bilen kjørt 20\enh{km} på 24\enh{min}, og må dermed har brutt fartsgrensen.

\begin{comment}
Se for deg at vi slipper en ball fra 60 meter over bakken, og lar den falle fritt nedover. Når vi slipper ballen, starter vi også en stoppeklokke. Antall meter $ h $ ballen er over bakken etter $ t $ sekunder kan da tilnærmes ved funksjonen\footnote{Funksjonsuttrykket kan faktoriseres, men vi har valgt å ikke gjøre det her.}
\[ h(t)=50-5t^2\quad,\quad x\in\left[0, \sqrt{10}\right] \]
\fig{ball}
Alle som har sett en ball falle, vet at den vil falle raskere og raskere fram til den treffer bakken. Ballen har altså \textit{ikke} konstant fart. Men det kan være nyttig å vite hva farten hadde vært \textit{hvis} den var konstant. La oss undersøke hva farten hadde vært hvis vi tenker oss at den var konstant mellom $ t=1 $ og $ t=3 $.
\fig{ballb}
Vi har at
\algv{
	h(1)&=50-5\cdot1^2 = 45 \\
	h(3)&=50-5\cdot3^2 = 5
}
Altså har ballen falt $ {(45-5)\enh{m}=40\enh{m}} $ på $ (3-1)\enh{s}=2\enh{s} $. Av \eqref{farteq} har vi da at
\[ \text{fart}=\frac{40\enh{m}}{2\enh{s}}=20\enh{m/s} \]
Som vi har sett på side ??, er dette \outl{gjennomsnittsfarten  til ballen} mellom tiden fra det 1. sekundet til det 3. sekundet. \newpage
Det er nå interessant å merke seg likheten mellom utrekningen av gjennomsnittsfarten og utrekningen for stigningstallet til linja mellom to punkt, som vi så på i \mb. Stigningstallet mellom punktene $ (1, h(1)) $ og $ (3, h(3)) $ er gitt som
\[ \frac{f(3)-f(1)}{3-1}=\frac{5-45}{2}=-20 \]
Gjennomsnittsfarten og stigningstallet har altså samme tallverdi, 20. Det som skiller verdiene, er at stigningstallet er negativt. Det forteller oss at ballen har gått i negativ retning, altså at ballen er fallende\footnote{Fordi vi har definert 'opp' som positiv retning.}. 	content...
\end{comment}
\end{document}

